\chapter{Introducción}

Este documento contiene el estudio de conexión simplificado del proyecto, el cual incluye los análisis de flujo de carga y cortocircuito para el sistema fotovoltaico de 141,6 kWp que se instalará en el supermercado MAS X MENOS sede Guarín, con conexión a la red de 13,8 kV propiedad de ESSA.

Este estudio debe presentarse ante el operador de red ESSA, y tiene como propósito determinar de qué manera se verá afectado el sistema eléctrico ante la entrada en operación de la nueva generación.

\section{Resumen Ejecutivo}

El supermercado MAS X MENOS sede Guarín contempla la instalación de un proyecto fotovoltaico de 141,6 kWp en sus instalaciones, cuya puesta en operación se encuentra en proceso de gestión.

La potencia instalada de generación en corriente alterna (AC) es de 120 kW, constituida por dos inversores SOLIS 60K-LV-5G de 60 kW cada uno, con tensión de salida de 220 V. Esta capacidad AC corresponde a una potencia instalada en corriente continua (DC) de 141,6 kWp, con una relación DC/AC de 1,18.

El proyecto se conectará a la red de ESSA a través de un transformador de 150 kVA, 13.200 V/220 V, en el nodo 3272966. Una vez entre en operación, el sistema suplirá durante algunos periodos del día la demanda total del cliente y exportará los excedentes de energía hacia el sistema de distribución local.

De los resultados obtenidos se destaca que, para la alternativa de conexión evaluada, el nuevo proyecto de generación no genera violaciones a los parámetros de operación normal de la red, y que las corrientes de cortocircuito no presentan un aumento significativo, debido a que la tecnología de generación empleada limita la corriente de inyección hacia el sistema.
